\chapter{Diskussion}
\label{ch:06_diskussion}

Dieses Kapitel führt die theoretischen Anforderungen und die praktischen Messergebnisse zusammen. Der Fokus liegt auf der Bewertung der Daten, der Vorstellung der heuristischen Entscheidungsmatrix als Kernbeitrag der Arbeit sowie einer kritischen Betrachtung der eigenen Arbeitsweise und der Erfahrungen bei der Entwicklung.


\section{Zusammenführung und Interpretation der Testergebnisse}
\label{sec:0601_interpretation_results}

Abgleich der gemessenen Kurven mit den zuvor aufgestellten Vermutungen aus den Tests. Beschreibung der genauen Belastungsgrenzen aller drei Render-Wege bei steigender Last durch Lichter, Geometriedichte und durcheinandergewürfelte Zeichenstile im Bild.


\section{Die heuristische Entscheidungsmatrix}
\label{sec:0602_heuristic_matrix}

Vorstellung des eigentlichen Ergebnisses der Arbeit in Form einer übersichtlichen Entscheidungsmatrix für Entwickler. Bereitstellung von konkreten Empfehlungen für den Aufbau einer Grafik-Engine basierend auf Geometriedichte, Lichtanzahl und gewünschten Zeichenstilen, wobei die Matrix in der praktischen Anwendung wie ein Kochbuch Schritt für Schritt durch die Architekturentscheidung führt.


\section{Kritische Reflexion und Limitationen}
\label{sec:0603_limitations}

Ehrliche Betrachtung der Schwachstellen und Grenzen des gewählten Versuchsaufbaus. Diskussion über bewusst weggelassene Faktoren wie andere Lichttypen, transparente Objekte oder die Ausführung auf mobilen Geräten und deren Auswirkung auf die Aussagekraft der Messergebnisse.


\section{Lessons Learned}
\label{sec:0604_lessons_learned_engineering}

Betrachtung von Nutzen und Aufwand aus Sicht der Softwareentwicklung abseits der reinen Performance-Messwerte. Vergleich zwischen der schwierigen Wartung von sehr großen Shader-Dateien und der besseren Modularität anderer Ansätze sowie Probleme beim fehlerfreien Einbau der Zeitmessung.